\section{Traditional Definition}
\label{sec:traditional-definition}

Given integers $\text{dividend}$ and $\text{divisor}$ where
$\text{divisor} \neq 0$, the division algorithm determines integers
$\text{quotient}$ and $\text{remainder}$ such that:

\begin{equation*}
\begin{aligned}
\forall\, \text{dividend},\text{divisor} \in \mathbb{Z},\;
\text{divisor} \neq 0,\;
\exists!\, \text{quotient},\text{remainder} &: \\
\text{dividend} &= \text{divisor} \cdot \text{quotient} + \text{remainder} \\
0 &\leq \text{remainder} < \lvert\text{divisor}\rvert \\
\text{dividend} \operatorname{div} \text{divisor} &:= \text{quotient} \\
\text{dividend} \operatorname{mod} \text{divisor} &:= \text{remainder}
\end{aligned}
\end{equation*}

The first two lines state the division relation and the canonical remainder
range. The final two lines introduce the operation notation: division returns
the quotient, and modulo returns the remainder.

\section{Recursive Definition}
\label{sec:recursive-definition}

We introduce a recursive definition of division and modulo because the proof
can be built from one invariant: shifting one unit of $b$ between quotient and
remainder preserves the represented dividend. In Section~5, that invariant
connects the recursive normal form back to the traditional division equation
from Section~3.

From here on, $a$ is the dividend, $b$ is the divisor, $q$ is the candidate
quotient, and $r$ is the candidate remainder. The shorter names keep the
recursive equations readable while preserving the same roles as the
traditional definition. We reserve $\operatorname{mod}$ for the modulo
operation itself. Throughout, $\mathbb{N}=\{0,1,2,\ldots\}$.

We define $\operatorname{DivMod}(a,b,q,r)$ such that:

\begin{equation*}
\begin{aligned}
\forall a,b,q,r \in \mathbb{Z} : b \neq 0,\; a = bq + r
\end{aligned}
\end{equation*}

The solved $\operatorname{DivMod}$ states are those where the remainder $r$
satisfies:

\begin{equation*}
\begin{cases}
0 \leq r < b & \text{if } b > 0, \\
0 \leq r < -b & \text{if } b < 0.
\end{cases}
\end{equation*}

\begin{equation*}
\begin{aligned}
\operatorname{DivMod.solve}(a,b,q,r) &:={}
\begin{cases}
\operatorname{DivMod}(a,b,q,r) & \text{if } 0 \leq r < \lvert b\rvert, \\
\operatorname{DivMod.solve}(a,b,q+\operatorname{sign}(b),r-\lvert b\rvert)
  & \text{if } r \geq \lvert b\rvert, \\
\operatorname{DivMod.solve}(a,b,q-\operatorname{sign}(b),r+\lvert b\rvert)
  & \text{if } r < 0.
\end{cases}
\end{aligned}
\end{equation*}

The recursive definition is implemented in
\href{https://github.com/thiagomata/prime-numbers/blob/modulo-article-v1.0.0/src/main/scala/v1/chapter2/div/DivMod.scala}{\texttt{DivMod.scala}}.
